\subsection{SEMI-SIMPLICITY OF REPRESENTATIONS}

Lemma 1. Let \(\mathfrak{g}\) be a semi-simple Lie algebra. The adjoint representation of \(\mathfrak{g}\) is semi-simple. Every ideal and every quotient algebra of \(\mathfrak{g}\) is semi-simple.

Let \(\mathfrak{a}\) be an ideal of \(\mathfrak{g}\) . The orthogonal \(\mathfrak{b}\) of \(\mathfrak{a}\) in \(\mathfrak{g}\) with respect to the Killing form is an ideal of \(\mathfrak{g}\) and \(\mathfrak{a} \cap \mathfrak{b}\) is a commutative ideal (§ 3, no. 6, Proposition 7) and hence zero. Hence \(\mathfrak{b}\) is supplementary to \(\mathfrak{a}\) in \(\mathfrak{g}\) . Moreover, as the Killing form of \(\mathfrak{g}\) is non-degenerate, so are its restrictions to \(\mathfrak{a}\) and \(\mathfrak{b}\) (Algebra, Chapter IX, § 4, no. 1, Corollary to Proposition 1) and hence \(\mathfrak{a}\) and \(\mathfrak{b}\) are semi-simple (no. 1, Theorem 1 and § 3, no. 6, Proposition 9).

Lemma 2. Let \(\mathfrak{g}\) be a Lie algebra. Then the following two conditions are equivalent:

\begin{enumerate}
\def\labelenumi{(\alph{enumi})}
\item
  All finite-dimensional linear representations of \(\mathfrak{g}\) are semi-simple.
\item
  Given a linear representation \(\rho\) of \(\mathfrak{g}\) on a finite-dimensional vector space V and a vector subspace W of codimension 1 such that \(\rho(x)(\mathbf{V}) \subset \mathbf{W}\) for all \(x \in \mathfrak{g}\) , there exists a supplementary line of W which is stable under \(\rho(\mathfrak{g})\) (and hence annihilated by \(\rho(\mathfrak{g})\) ).
\end{enumerate}

Clearly (a) implies (b). Suppose that (b) holds. Let \(\sigma\) be a finite-dimensional representation of \(\mathfrak{g}\) on a vector space M and N a vector subspace which is stable under \(\sigma(\mathfrak{g})\) . Let \(\mu\) be the representation of \(\mathfrak{g}\) on \(\mathcal{L}(\mathrm{M})\) canonically derived from \(\sigma\) (§ 3, no. 3): recall that \(\mu(x) = \mathrm{ad}_{\mathcal{L}(\mathrm{M})}\sigma(x)\) . Let V (resp. W) be the subspace of \(\mathcal{L}(\mathrm{M})\) consisting of the linear mappings of M into N whose restriction to N is a homothety (resp. zero); then W is of codimension 1 in V and \(\mu(x)(\mathrm{V})\subset\mathrm{W}\) for all \(x\in\mathfrak{g}\) . By condition (b), there exists \(u\in\mathrm{V}\) which is annihilated by \(\mu(x)\) for all \(x\in\mathfrak{g}\) and whose restriction to N is a non-zero homothety. By multiplying \(u\) by a suitable scalar, it can be assumed that \(u\) is a projector of M onto N. To say that \(\mu(x).u=0\) means that \(u\) is permutable with \(\sigma(x)\) . Hence the kernel of \(u\) is a supplement of N in M which is stable under \(\sigma(x)\) for all \(x\in\mathfrak{g}\) . Hence \(\sigma\) is semi-simple.

Lemma 3. Let g be a semi-simple Lie algebra, \(\rho\) a linear representation of \(g\) on a finite-dimensional vector space \(V\) and \(W\) a subspace of \(V\) of codimension 1 such that \(\rho(x)(V) \subset W\) for all \(x \in g\) . Then there exists a supplementary line of \(W\) which is stable under \(\rho(g)\) .

For all \(x \in \mathfrak{g}\) let \(\sigma(x)\) be the restriction of \(\rho(x)\) to W. Suppose first that \(\sigma\) is simple. If \(\sigma = 0\) , then \(\rho(x)\rho(y) = 0\) for all \(x,y\) in \(\mathfrak{g}\) , hence \(\rho(\mathfrak{g}) = \rho(\mathcal{D}\mathfrak{g}) = \{0\}\) and our assertion is obvious. If \(\sigma \neq 0\) , let \(\mathfrak{n}\) be the kernel of \(\sigma\) and let \(\mathfrak{m}\) be a supplementary ideal of \(\mathfrak{n}\) in \(\mathfrak{g}\) (Lemma 1); then \(\mathfrak{m} \neq \{0\}\) and the restriction of \(\sigma\) to \(\mathfrak{m}\) is faithful; the restriction to \(\mathfrak{m}\) of the bilinear form associated with \(\sigma\) is non-degenerate (Proposition 1) and hence the Casimir element \(c\) associated with \(\mathfrak{m}\) and \(\sigma\) can be formed. By Proposition 12 of §3, no. 7, \(\sigma(c)\) is an automorphism of W. On the other hand, \(\rho(c)(\mathrm{V}) \subset \mathrm{W}\) . Hence the kernel Z of \(\rho(c)\) is a supplementary line of W; since \(c\) belongs to the centre of the enveloping algebra of \(\mathfrak{g}\) , \(\rho(c)\) is permutable with \(\rho(x)\) for all \(x \in \mathfrak{g}\) and hence Z is stable under \(\rho(\mathfrak{g})\) .

In the general case we argue by induction on the dimension of V. Let T be a minimal non-zero stable subspace of W. Let \(\rho'\) be the quotient representation on \(V' = V/T\) . Then, for all \(x \in g\) , \(\rho'(x)(V') \subset W'\) , where \(W' = W/T\) is of codimension 1 in \(V'\) . By the induction hypothesis there exists a line \(Z'\) which is supplementary to \(W'\) and stable under \(\rho'(g)\) . Its inverse image Z in V is stable under \(\rho(g)\) , contains T as subspace of codimension 1, \(Z \cap W = T\) , and hence \(\rho(x)(Z) \subset T\) for all \(x \in g\) . By what was proved above, there exists a supplementary line of T in Z which is stable under \(\rho(g)\) ; this line is supplementary to W in V, which completes the proof.

THEOREM 2 (H. Weyl). Every finite-dimensional linear representation of a semi-simple Lie algebra is completely reducible.

This follows from Lemmas 2 and 3.

DEFINITION 2. A Lie algebra g is called simple if the only ideals of g are \{0\} and g and if further g is non-commutative.

A simple Lie algebra is semi-simple. The algebra \(\{0\}\) is not simple.

PROPOSITION 2. For a Lie algebra g to be semi-simple, it is necessary and sufficient that it be a product of simple algebras.

The condition is sufficient (no. 1, Remark 3). Conversely, suppose that g is semi-simple. Since the adjoint representation of g is semi-simple, g is the direct sum of minimal non-zero ideals \(a_{1}, \ldots, a_{m}\) . Then g is identified with the product algebra of the \(a_{i}\) (§ 1, no. 1). Every ideal of \(a_{i}\) is then an ideal of g and hence zero or equal to \(a_{i}\) . On the other hand \(a_{i}\) is non-commutative. Hence the \(a_{i}\) are simple Lie algebras.

COROLLARY 1. A semi-simple Lie algebra is the product of its simple ideals \(\mathfrak{g}_{\mathrm{i}}\) . Every ideal of \(\mathfrak{g}\) is the product of certain of the \(\mathfrak{g}_{\mathrm{i}}\) .

\(g = a_{1} \times \cdots \times a_{m}\) , where the \(a_{i}\) are simple. As the centre of \(a_{i}\) is zero, the centralizer of \(a_{i}\) in g is the product of the \(a_{j}\) for \(j \neq i\) . Then let a be an ideal of g. If it does not contain \(a_{i}\) , then \(a \cap a_{i} = \{0\}\) , hence \([a, a_{i}] = \{0\}\) and a is contained in the product of the \(a_{j}\) for \(j \neq i\) . It follows that a is the product of certain of the \(a_{i}\) . Hence the simple ideals of g are precisely the \(a_{i}\) .

The simple ideals of a semi-simple Lie algebra are called the simple components of g.

COROLLARY 2. Let \(\mathfrak{g},\mathfrak{g}'\) be two Lie algebras, \(\mathfrak{r}\) and \(\mathfrak{r}'\) their radicals and \(f\) a homomorphism of \(\mathfrak{g}\) onto \(\mathfrak{g}'\) . Then \(\mathfrak{r}' = f(\mathfrak{r})\) .

As \(f(\mathfrak{r})\) is solvable, \(f(\mathfrak{r}) \subset \mathfrak{r}'\) . On the other hand, \(\mathfrak{g} / \mathfrak{r}\) is semi-simple (§ 5, no. 2, Proposition 3), hence \(\mathfrak{g}' / f(\mathfrak{r})\) , which is isomorphic to a quotient of \(\mathfrak{g} / \mathfrak{r}\) , is semi-simple (Lemma 1) and hence \(f(\mathfrak{r}) \supset \mathfrak{r}'\) (§ 5, no. 2, Proposition 3).

Remarks. (1) Theorem 2 admits a converse: if every finite-dimensional representation of g is semi-simple, g is semi-simple. For since the adjoint representation is semi-simple, every ideal of g admits a supplementary ideal and hence can be considered as a quotient of g. If g is not semi-simple then g admits a non-zero commutative quotient and therefore a quotient of dimension 1. Now the Lie algebra K of dimension 1 admits non-semi-simple representations, for example

\[
\lambda \mapsto \left( \begin{array}{c c} 0 & 0 \\ \lambda & 0 \end{array} \right).
\]

\begin{enumerate}
\def\labelenumi{(\arabic{enumi})}
\setcounter{enumi}{1}
\tightlist
\item
  Let g be a Lie algebra over K and σ a representation of g on a vector space M. On the other hand let \(f\) be a K-linear mapping of \(g\) into \(M\) such that:
\end{enumerate}

\[
f ([ x, y ]) = \sigma (x). f (y) - \sigma (y). f (x)\tag{1}
\]

for all \(x, y\) in \(\mathfrak{g}\) . By § 1, no. 8, Example 2, being given \(\sigma\) and \(f\) is equivalent to being given a homomorphism \(x \mapsto (f(x), \sigma(x))\) of \(\mathfrak{g}\) into \(\mathfrak{af}(\mathbf{M})\) . On the other hand we have seen (loc. cit.) that the element \((f(x), \sigma(x))\) of \(\mathfrak{af}(\mathbf{M})\) is canonically identified with the element \(\rho(x)\) of \(\mathfrak{gl}(\mathbf{N})\) (where \(\mathrm{N} = \mathrm{M} \times \mathrm{K}\) ) which induces \(\sigma(x)\) on \(\mathrm{M}\) and maps the element (0, 1) of \(\mathrm{N}\) to \(f(x)\) . And \(\rho\) is then a representation of \(\mathfrak{g}\) on \(\mathrm{N}\) such that \(\rho(x)(\mathrm{N}) \subset \mathrm{M}\) for all \(x \in \mathfrak{g}\) .

Then, if g is semi-simple, there exists (Lemma 3) a line Z which is supplementary to M in N and annihilated by \(\rho(\mathfrak{g})\) . In other words, there exists an element \(m_{0} \in M\) such that \((-m_{0}, 1) \in N\) is annihilated by \(\rho(x)\) , that is such that

\[
f (x) = \sigma (x). m _ {0}\tag{2}
\]

for all \(x \in g\) .

*Suppose that K = R. Let G be a connected Lie group with Lie algebra g. Consider an analytic homomorphism \(\phi\) of G into the affine group A of M corresponding to a homomorphism \(x \mapsto (f(x), \sigma(x))\) of g into \(\mathfrak{af}(M)\) . The above results can be interpreted by saying that if g is semi-simple \(\phi(G)\) leaves a point of M fixed. For let H be the set of elements of GL(N) which leave stable all the linear varieties of N parallel to M. There exists (§ 1, no. 8, Example 2) a canonical isomorphism \(\psi\) of A onto H. Let Z be a supplementary line of M in N. To say that \(\rho(\mathfrak{g})\) annihilates Z amounts to saying that \((\psi \circ \phi)(G)\) leaves the points of Z fixed and hence (taking account of the definition of \(\psi\) ) that \(\phi(G)\) leaves fixed the projection onto M of the point of intersection of Z and M × \{1\}.

